\subsection{Preliminaries} \label{sec:prelim}
We use $[n]$ to denote the set $1, \ldots, n$.  We write $x \gets X$ to indicate that the value $x$ is sampled uniformly at random from the set $X$.  We use $\secparam$ to denote a statistical security parameter and $\csecparam$ to denote a computational security parameter.  We also assume the existence of a hash function $H:\{0,1\}^* \rightarrow \{0,1\}^*$ which is modeled as a random oracle.  We let $\poly(\cdot)$ denote a polynomial function and $\negl(\cdot)$ denote a negligible function.  
%Some important variable names and notation are given in Figure~\ref{fig:notation}.

% \begin{figure}
% \begin{tabular}{|c|c|}
% \hline
% Symbol     &  Meaning\\
% \hline
% %$A$ & Sender\\
% %$B$ & Receiver \\
% %$C$ & Complainer \\
% %$O$ & Originator (used sometimes) \\
% &\textbf{FACTS variables}\\
% $S$ & Server \\
% $c$ & Number of users \\
% $x$ & Message or item \\
% $\secparam$ & statistical security parameter \\
% $\csecparam$ & computational security parameter\\
% &\textbf{CCBF variables}\\
% %$h$ & hash of message \\
% %$e$ & encryption of user ID\\
% $T$ & bit vector stored by the server \\
% $s$ & size of T\\
% $n$ & total number of complaints per epoch \\
% $t$ & (explicit) threshold \\
% $u$ & size of user set \\
% $v$ & size of item set\\
% %$\mathsf{tag}$ & tag\\
% %$\tau$ & tipping point / implicit threshold \\
% %$r$ & salt for hash of message \\
% %$m$ & current total number of complaints \\
% %$\ell$ & upper bound of lists \\
% %$w$ & current \# of free slots for a given message \\
% %$C$ & user id \\
% %$x$ & item / message \\
% $U_C$ & set of bit indices for user $C$ \\
% $V_x$ & set of bit indices for item/message $x$ \\
% %$i$ & arbitrary index value used in some proof \\
% %$p$ & probability \\
% %$H$ & Hash function \\
% %$\s$ & Signature \\
% %$L$ & daily limit on number of complaints per party\\
% \hline
% \end{tabular}
% \caption{Notation}\label{fig:notation}
% \end{figure}

